% !PRINTED 119
\weberpage{119}

This is a defect of \gesperrt{Cardano's formula}, which accordingly by itself does not yet suffice to give the roots of the cubic equation unambiguously. \gesperrt{Cayley} remedied this defect by the following representation. Define two new quantities $\xi,\eta$ by the equations
\begin{equation}
u=\xi^2\eta,
\qquad
v=\xi\eta^2,
\tag{2}
\end{equation}
from which, by multiplication and use of (1), one obtains
\begin{equation}
\xi\eta=\sqrt[3]{-\frac a3}.
\tag{3}
\end{equation}
Consequently, if this value is inserted in (2) and the expressions \S\,35, (8) for $u,v$ are substituted, then
\begin{equation}
\xi=\sqrt[3]{\frac{3b}{2a}+\sqrt{\frac{9b^2}{4a^2}+\frac a3}},\qquad
\eta=\sqrt[3]{\frac{3b}{2a}-\sqrt{\frac{9b^2}{4a^2}+\frac a3}}.
\tag{4}
\end{equation}
Thus, if one now puts the root $y$ of the cubic equation in the form
\begin{equation}
y=\xi\eta(\xi+\eta),
\tag{5}
\end{equation}
one obtains an expression which, when $\xi$ is replaced, independently, by $\xi,\varepsilon\xi,\varepsilon^2\xi$ and $\eta$ by $\eta,\varepsilon\eta,\varepsilon^2\eta$, assumes not nine but \gesperrt{only the three values}\footnotemark
\[
\begin{array}{c}
\xi^2\eta+\xi\eta^2,\\
\varepsilon\xi^2\eta+\varepsilon^2\xi\eta^2,\\
\varepsilon^2\xi^2\eta+\varepsilon\xi\eta^2
\end{array}
\]
\footnotetext{\gesperrt{Cayley}, Phil. Mag. vol. XXI, 1861. Collectet mathematical papers vol. V, No. 310.}% [sic] The registered source prints ‘Collectet’; preserved as a source-original bibliographic anomaly. The callout is attached to the translated phrase introducing the three displayed values so it cannot strand after the display.



\begin{center}
{\large \S.\ 37.}

\vspace{0.5em}
{\large\gesperrt{The Biquadratic Equation.}}
\end{center}

A similar path to that used in solving the cubic equation by \gesperrt{Cardano's formula} may be taken for the solution of the equation of the fourth degree
\begin{equation}
 y^4+ay^2+by+c=0.
\tag{1}
\end{equation}
We put
%CONT>>
